\scanpage{146}
\margin{\textcircled{b}}
$L_\bullet$ un complexe pseudo-cohérent de Tor-dimension
finie. Alors
\diagram{\begin{tikzcd}[column sep=large,row sep=large]
A_X\arrow[r,"\varepsilon"]\arrow[dr,"\unsure{\chi|\gamma}_{L_\bullet}\deleted{\unsure{*}}"']&\mathbb R\mathcal Hom(L_\bullet,L_\bullet)\arrow[d,"\tau"]&L_\bullet^\vee\overset{\mathbb L}{\otimes}L_\bullet\arrow[l,"\simeq"']\\
&A_X&
\end{tikzcd}}
Cet accouplement induit une forme linéaire
\[
\mathcal H^0(\mathbb R\mathcal Hom(L_\bullet,L_\bullet))\longrightarrow A_X
\]
qui s'explicite, lorsque $L_\bullet$ est \deleted{plat}\inserted{loc. libre de t.f., \deleted{et borné},} en chaque
degré, comme
\[
\varphi\longmapsto\sum(-1)^i\operatorname{Tr}\varphi_i^{\deleted{\unsure{0}}}
\]
si $\varphi_\bullet=(\unsure{\varphi_i|\varphi})$ est \deleted{\unsure{un hom}} un hom de
complexes $L_\bullet\longrightarrow L_\bullet$.
\deleted{[Regardons en effet}
\deleted{Si $\varphi_i\in L_i^\vee\otimes L_i$, sa \illegible{1} \unsure{lin.} associée}

Supposons maintenant que $L_\bullet,M_\bullet$
soient tous deux de Tor-dim finie
et pseudo-cohérents, considérons les accouplements
naturels,
\[
\begin{gathered}
(u,v)\longmapsto vu\\
\mathbb R\mathcal Hom^\bullet(L_\bullet,M_\bullet)\times\mathbb R\mathcal Hom(M_\bullet,L_\bullet)\xrightarrow{\varphi_{L_\bullet,M_\bullet}}\mathbb R\mathcal Hom^\bullet(L_\bullet,L_\bullet)\\
\mathbb R\mathcal Hom^\bullet(M_\bullet,L_\bullet)\times\mathbb R\mathcal Hom^\bullet(L_\bullet,L_\bullet)\longrightarrow\mathbb R\mathcal Hom^\bullet(M_\bullet,M_\bullet)\\
(v,u)\longmapsto uv
\end{gathered}
\]
alors \deleted{les \illegible{1} \unsure{préc.}} on a
\[
\tau_{L_\bullet}\varphi_{L_\bullet,M_\bullet}=\tau_{M_\bullet}\varphi_{M_\bullet,L_\bullet}\operatorname{sym}
\]
\mspage{145}
